% ============================================================
% Estudo de viabilidade (congelado) do projeto
% "Da rede lexical à rede contextual" -- MO438 Redes Complexas, Unicamp.
%
% Compilação:  cd report/estudo-de-viabilidade && latexmk -pdf estudo-de-viabilidade.tex
%
% Os números deste documento são de planejamento: foram copiados das saídas dos scripts em
% scripts/feasibility/ (sim_<modelo>.txt, hyb_<modelo>.txt) e não mudam com o pipeline.
% As figuras (figuras/viabilidade_*.pdf) vêm de scripts/feasibility/plot_feasibility.py.
% ============================================================
\documentclass[11pt,a4paper]{article}

\usepackage[utf8]{inputenc}
\usepackage[T1]{fontenc}
\usepackage[brazil]{babel}
\usepackage{lmodern}
\usepackage{microtype}
\usepackage{amsmath,amssymb}
\usepackage{graphicx}
\usepackage{booktabs}
\usepackage{geometry}
\usepackage{caption}
\usepackage{hyperref}

% O mesmo layout do relatório técnico: as figuras foram desenhadas para essa largura.
\geometry{top=2.0cm, bottom=2.0cm, left=2.0cm, right=2.0cm}
\setlength{\parindent}{0pt}
\setlength{\parskip}{4pt}
\setlength{\emergencystretch}{2.5em}
\captionsetup{font=small, labelfont=bf, margin=0.5cm}
\captionsetup[table]{position=top, skip=4pt}
\hypersetup{
    colorlinks=true,
    linkcolor=black,
    urlcolor=blue,
    pdftitle={Da rede lexical à rede contextual: estudo de viabilidade},
    pdfauthor={Juliana Midlej},
}

\newcommand{\tok}[1]{\texttt{#1}}              % um token, como aparece no texto
\newcommand{\esp}{\textvisiblespace}           % espaço inicial de um token
\newcommand{\cod}[1]{\texttt{#1}}              % código, arquivos, comandos
% \figura{arquivo sem extensão}{legenda}{rótulo}
\newcommand{\figura}[3]{%
    \begin{figure}[htbp]
        \centering
        \includegraphics{figuras/#1.pdf}
        \caption{#2}\label{#3}
    \end{figure}}
\newcommand{\repo}{https://github.com/annwith/gender-representation-networks}

% ============================================================
\begin{document}

\begin{center}
    {\large\bfseries MO438 -- Redes Complexas (Unicamp)}\\[4pt]
    {\LARGE\bfseries Da rede lexical à rede contextual}\\[3pt]
    {\large Estudo de viabilidade}\\[8pt]
    Juliana Midlej
\end{center}

\paragraph{O que é este documento.} O projeto representa a passagem das representações
lexicais às contextuais do Qwen3-4B-Base como uma sequência de redes $k$-NN, em que os vértices
são \emph{ocorrências} de tokens em parágrafos da Wikipédia em português. Uma das perguntas, a
P4, compara as ocorrências de uma mesma palavra em temas diferentes (por exemplo,
\emph{banco} em Economia e em Computação), e por isso exige palavras repetidas o bastante em
vários temas. Este estudo foi feito antes do pipeline, para dimensionar a amostra e escolher o
modelo. Ele está congelado: os números são de planejamento, vêm de um corpus simplificado e não
mudam com as execuções. O método e os resultados do experimento estão no relatório técnico,
no repositório (\url{\repo}).

\textbf{Notação.} $f_t$ é o número de vértices do tipo (token) $t$ na amostra, e $f_{\max}$ é o
limite de vértices por tipo.

Antes de implementar o pipeline, um estudo rápido (scripts em \cod{scripts/feasibility/})
respondeu a três perguntas: (i) a meta de cerca de 15 mil vértices é alcançável com parágrafos
inteiros? (ii) Parágrafos inteiros deixam palavras repetidas suficientes, em temas diferentes,
para a P4? (iii) Que tokenizer, e portanto que modelo, usar?

\paragraph{Dados do estudo.} O \cod{fetch\_wiki.py} baixou uma versão simplificada do corpus:
os mesmos oito temas, busca até a profundidade 2, 120 artigos por tema, limpeza só com
\cod{strip\_code} e os mesmos filtros de parágrafo. Resultaram \textbf{958 artigos} e
\textbf{7.876 parágrafos} (Biologia 965, Computação 1.034, Economia 1.147, Esporte 836, Física
948, Geografia 785, Música 1.032, Política 1.129), com \textbf{1,28 milhão de tokens} do
Qwen3 (1.276.475) e cerca de 23 mil tipos distintos. Três tokenizers foram comparados: o do
Qwen3-4B (idêntico ao do Qwen3-4B-Base, 151.669 ids), o do Tucano-2b4 (treinado em português) e
o do GPT-2 small português (o modelo do piloto).

\paragraph{Amostragem natural.} O \cod{simulate\_sample.py} simula a amostragem mais simples:
parágrafos inteiros, em rodízio entre os temas, com todos os tokens como vértices e um limite
$f_{\max}$ por tipo, até chegar a 15 mil vértices. Também simula uma variante
``direcionada'', que começa pelos parágrafos que contêm as palavras-alvo (até 12 por palavra e
tema) e depois segue a ordem natural. As saídas estão em \cod{sim\_<modelo>.txt}.

\paragraph{Desenho híbrido.} O \cod{sim\_hybrid.py} simula a alternativa adotada: um
\emph{núcleo} denso de parágrafos inteiros mais estratos esparsos de ocorrências escolhidas
(palavras-alvo, controles e palavras de conteúdo presentes em vários temas), em que só as
posições sorteadas viram vértices e o resto do parágrafo é contexto. As saídas estão em
\cod{hyb\_<modelo>.txt}.

Nas contagens a seguir, \emph{palavra de conteúdo} é um tipo cuja categoria mais comum é
palavra inteira, só com letras, com pelo menos 3 letras e fora de uma lista de palavras
funcionais. O critério da P4 é ter $f_t\ge10$ na amostra e pelo menos 3 ocorrências em pelo
menos 3 temas.

\begin{table}[htbp]
\centering
\caption{Estudo de viabilidade: tokenizers e amostragem natural até 15 mil vértices com
$f_{\max}=50$. Fontes: linha \cod{pool:} e bloco \cod{f\_max=50, amostragem natural} de
\cod{scripts/feasibility/sim\_<modelo>.txt}.}\label{tab:viab-natural}
\small
\begin{tabular}{@{}l r r r@{}}
\toprule
& \textbf{Qwen3-4B} & \textbf{Tucano-2b4} & \textbf{GPT-2 small PT} \\
\midrule
Tokens no \emph{pool} de 7.876 parágrafos & 1.276.475 & 953.290 & 939.895 \\
Tokens por palavra & 1,78 & 1,33 & 1,31 \\
Palavras-alvo candidatas que são um token só & 12 de 16 & 16 de 16 & 16 de 16 \\
\midrule
Parágrafos para chegar a 15 mil vértices & 101 & 134 & 160 \\
Tipos distintos na amostra & 3.693 & 6.144 & 6.213 \\
Vértices com $f_t\le2$ & 19\% & 40\% & 39\% \\
Vértices que são palavras inteiras & 28\% & 54\% & 72\% \\
Tipos com $f_t\ge10$ & 293 & 183 & 170 \\
Conteúdo com $f_t\ge10$ e $\ge3$ oc.\ em $\ge3$ temas & 4 & 5 & 9 \\
Alvos com $\ge10$ oc.\ em $\ge2$ temas & 0 & 0 & 0 \\
\midrule
\multicolumn{4}{@{}l}{\emph{Variante direcionada (alvos primeiro), $f_{\max}=50$}} \\
Conteúdo com $f_t\ge10$ e $\ge3$ oc.\ em $\ge3$ temas & 10 & 10 & 21 \\
Alvos com $\ge10$ oc.\ em $\ge2$ temas & 2 & 1 & 2 \\
\bottomrule
\end{tabular}
\end{table}

\figura{viabilidade_tokenizers}
    {Comparação dos tokenizers no \emph{pool} do estudo de viabilidade: tokens por palavra e
    fração dos vértices que são palavras inteiras na amostragem natural ($f_{\max}=50$). O
    Qwen3 (em destaque) é o tokenizer do modelo escolhido. Fonte:
    \cod{scripts/feasibility/sim\_<modelo>.txt}.}
    {fig:viab-tokenizers}

A Tabela~\ref{tab:viab-natural} e a Figura~\ref{fig:viab-tokenizers} mostram três coisas.

\textbf{A meta de vértices é fácil.} Com qualquer tokenizer, 15 mil vértices saem de 100 a 160
parágrafos (de 18 a 21 mil tokens), uma fração pequena do \emph{pool}. Com $f_{\max}=20$ bastam
de 127 a 177 parágrafos.

\textbf{Parágrafos inteiros não servem para a P4.} Mesmo com 15 mil vértices, só de 4 a 9
palavras de conteúdo têm $f_t\ge10$ com ocorrências espalhadas por 3 temas, e \emph{nenhuma}
palavra-alvo chega a 10 ocorrências em 2 temas. O motivo é a lei de Zipf: num trecho pequeno de
texto, só palavras funcionais e pontuação se repetem muito; palavras de conteúdo aparecem uma
ou duas vezes. Começar pelos parágrafos que contêm as palavras-alvo (variante direcionada)
melhora pouco (de 10 a 21 palavras de conteúdo, 1 ou 2 alvos).

\textbf{O Qwen3 fragmenta muito o português.} São 1,78 tokens por palavra, contra 1,33 do
Tucano e 1,31 do GPT-2 português, e só 28\% dos vértices são palavras inteiras (54\% e 72\% nos
outros). Quatro das 16 palavras-alvo candidatas não são um token só no Qwen3 (\emph{corrente},
\emph{onda}, \emph{célula} e \emph{núcleo}; por exemplo, ``\esp célula'' vira
\tok{\esp cél}${}+{}$\tok{ula}) e saíram da lista, que ficou com 12. Essa foi a principal
desvantagem pesada contra o Qwen3 na escolha do modelo (relatório técnico, seção ``Escolha do modelo'').

\begin{table}[htbp]
\centering
\caption{Estudo de viabilidade: desenho híbrido com 15 mil vértices e $f_{\max}=50$
(núcleo de parágrafos inteiros, 12 ou 16 alvos, 6 controles e 150 palavras multitema com 20
ocorrências cada). Fonte: \cod{scripts/feasibility/hyb\_<modelo>.txt}.}\label{tab:viab-hibrido}
\small
\begin{tabular}{@{}l r r r@{}}
\toprule
& \textbf{Qwen3-4B} & \textbf{Tucano-2b4} & \textbf{GPT-2 small PT} \\
\midrule
Vértices & 15.104 & 15.002 & 15.099 \\
\quad núcleo & 11.251 & 10.954 & 11.045 \\
\quad alvos & 553 & 748 & 754 \\
\quad controles & 300 & 300 & 300 \\
\quad multitema & 3.000 & 3.000 & 3.000 \\
Parágrafos inteiros no núcleo & 78 & 93 & 107 \\
Vértices por parágrafo do núcleo (média) & 144 & 118 & 103 \\
Parágrafos tocados no total & 2.894 & 2.943 & 3.017 \\
Tokens a processar no modelo & 329.555 & 243.959 & 248.579 \\
Tipos distintos & 3.257 & 4.904 & 4.856 \\
Vértices com $f_t\le2$ & 17\% & 32\% & 30\% \\
Tipos com $f_t\ge10$ & 354 & 274 & 286 \\
Conteúdo com $f_t\ge10$ e $\ge3$ oc.\ em $\ge3$ temas & 169 & 172 & 173 \\
Alvos com $\ge10$ oc.\ em $\ge2$ temas & 3 de 12 & 7 de 16 & 7 de 16 \\
Alvos com $\ge5$ oc.\ em $\ge3$ temas & 11 de 12 & 14 de 16 & 14 de 16 \\
\bottomrule
\end{tabular}
\end{table}

\figura{viabilidade_amostragem}
    {Amostragem natural (parágrafos inteiros, $f_{\max}=50$) contra o desenho híbrido, nas duas
    contagens que decidem a P4: palavras de conteúdo repetidas em vários temas e palavras-alvo
    com ocorrências suficientes em dois temas. Fontes:
    \cod{scripts/feasibility/sim\_<modelo>.txt} e \cod{hyb\_<modelo>.txt}.}
    {fig:viab-amostragem}

\textbf{O desenho híbrido resolve.} Com o mesmo número de vértices, o núcleo continua com cerca
de 11 mil vértices em 78 a 107 parágrafos inteiros (o bastante para a P3 e para as métricas
globais), e o número de palavras de conteúdo com $f_t\ge10$ espalhadas por 3 temas sobe de
4--9 para cerca de 170 (Tabela~\ref{tab:viab-hibrido} e Figura~\ref{fig:viab-amostragem}). O
custo é processar mais texto no modelo: cerca de 330 mil tokens com o Qwen3, porque cada
ocorrência esparsa exige o parágrafo até a posição dela.

Na simulação, os alvos ainda ficaram fracos (3 de 12 com $\ge10$ ocorrências em 2 temas), mas
por uma razão corrigível: as 50 ocorrências de cada alvo foram sorteadas em rodízio entre os
\emph{oito} temas, e não concentradas nos temas em que a palavra tem sentidos diferentes. A
Tabela~\ref{tab:viab-pool} mostra que o \emph{pool} tem ocorrências suficientes para cotas por
\emph{tema de sentido} na maioria dos casos, e onde faltam (\emph{banco} em Geografia,
\emph{órgão} em Biologia, \emph{nota} e \emph{carga} em Economia, \emph{planta} em Economia,
\emph{matriz} em Computação), entram as categorias extras do corpus final (relatório técnico, seção ``Corpus''). A
simulação também tinha dois defeitos de desenho, corrigidos no desenho final
(relatório técnico, problema ``Limite por ordem de chegada e artigo igual a parágrafo
na simulação''): o limite por tipo do núcleo era aplicado por ordem de
chegada, e cada parágrafo do núcleo vinha de um artigo diferente.

\begin{table}[htbp]
\centering
\caption{Ocorrências das palavras-alvo e de controle no \emph{pool} inteiro do estudo de
viabilidade (tokenizer do Qwen3, sem limite por tipo), por tema. Em negrito, os temas de
sentido usados nas cotas do desenho final (relatório técnico, tabela de cotas da seção ``Amostra híbrida''). Fonte: bloco \cod{pool
inteiro} de \cod{scripts/feasibility/sim\_Qwen3-4B.txt}.}\label{tab:viab-pool}
\footnotesize
\setlength{\tabcolsep}{4.5pt}
\begin{tabular}{@{}l r r r r r r r r r@{}}
\toprule
& \textbf{Total} & \textbf{Bio} & \textbf{Comp} & \textbf{Econ} & \textbf{Esp} & \textbf{Fís}
& \textbf{Geo} & \textbf{Mús} & \textbf{Pol} \\
\midrule
banco & 66 & 9 & \textbf{19} & \textbf{24} & 6 & 0 & \textbf{4} & 3 & 1 \\
campo & 335 & 34 & 26 & 22 & \textbf{74} & \textbf{82} & \textbf{72} & 11 & 14 \\
órgão & 57 & \textbf{6} & 1 & 6 & 9 & 0 & 0 & \textbf{25} & \textbf{10} \\
estado & 282 & 24 & 39 & 28 & 13 & \textbf{84} & 19 & 20 & \textbf{55} \\
carga & 53 & 9 & 4 & \textbf{6} & 0 & \textbf{22} & 0 & 3 & 9 \\
rede & 133 & 4 & \textbf{71} & 15 & 2 & 11 & \textbf{15} & 3 & 12 \\
nota & 41 & 2 & 5 & \textbf{7} & 0 & 2 & 0 & \textbf{24} & 1 \\
capital & 271 & 1 & 0 & \textbf{201} & 18 & 1 & 9 & 6 & \textbf{35} \\
família & 70 & \textbf{10} & 10 & 4 & 4 & 13 & 2 & 13 & \textbf{14} \\
massa & 140 & 2 & 3 & 10 & 3 & \textbf{59} & \textbf{39} & 14 & 10 \\
matriz & 27 & \textbf{10} & \textbf{3} & 2 & 0 & 5 & 0 & 6 & 1 \\
planta & 35 & \textbf{25} & 0 & \textbf{0} & 0 & 0 & 9 & 0 & 1 \\
\midrule
ano & 310 & 13 & 15 & \textbf{36} & \textbf{82} & 35 & 30 & \textbf{65} & 34 \\
século & 356 & 29 & 12 & 31 & 11 & \textbf{86} & \textbf{66} & 59 & \textbf{62} \\
grande & 527 & \textbf{58} & 48 & 67 & 66 & \textbf{88} & 49 & 77 & 74 \\
primeiro & 510 & 24 & \textbf{52} & 39 & \textbf{134} & 63 & 48 & 92 & 58 \\
importante & 247 & \textbf{53} & 20 & \textbf{47} & 9 & 24 & 23 & 36 & 35 \\
água & 249 & \textbf{102} & 0 & 24 & 2 & 47 & \textbf{65} & 1 & 8 \\
\bottomrule
\end{tabular}
\end{table}

Os números do estudo são aproximados por construção: o corpus do estudo não removia páginas
multitema nem biografias, a limpeza ainda deixava restos de referências e as contagens
ignoram os limites por parágrafo e por artigo do desenho final. Eles serviram para escolher o
desenho e o modelo; a composição real da amostra está na seção ``Amostra híbrida'' do relatório técnico.

\end{document}
